% ARGUES: this is the experimental complement to the statistical account of
%         hallucination, and it is not scooped by the concurrent solution-hacking work.
% DRAWS FROM: docs/related_work_solution_hacking.md, PRE_REGISTRATION_OSF.md §1.
% STATUS: drafted 2026-08-12. Every citation in refs.bib is still unverified.

\section{Related work}

\subsection{Hallucination as a statistical equilibrium}

\citet{kalai2025} give a statistical account of why hallucination persists.
Standard training and evaluation reward confident guessing over calibrated
abstention, so producing an answer is optimal under a misspecified scoring rule
even when the model has no basis for one. Hallucination is the equilibrium of
that rule rather than a defect that more data removes.

Our work is the experimental complement to that argument, not a correction of
it. The difference is in what each can establish. \citet{kalai2025} name both
capability (the model cannot retrieve the fact) and incentive (the model is
rewarded for not abstaining) as contributors, but the analysis is of the
training objective and does not separate the two empirically. The
obscure-but-real cell makes that separation a controllable contrast: capability
is held near-constant against the unanswerable cell while what an honest answer
looks like changes. We do not attribute to \citet{kalai2025} a prediction in the
form of a manipulable cell, because they make none.

The other differences follow from that. Their analysis is of training time and
manipulates nothing inside a fixed policy, whereas E2 does. And where they
predict the equilibrium, we measure whether four separate things move in the
direction the reward-hacking account requires: the post-training gradient, the
prompt lever, the internal state, and a plausible surface proxy.

\subsection{Reward hacking, and what this paper does not claim}

We adopt the definition of reward hacking in \citet{skalse2022}: a policy scores
well on a proxy while doing worse on the objective the proxy stands in for.
Citation fabrication fits that shape unusually cleanly, because the proxy
(an answer that looks sourced) and the objective (an answer that is sourced) come
apart only when someone checks, and checking is the expensive step that
preference pipelines skip.

\citet{hubinger2019} on deceptive alignment and \citet{park2024} on strategic
deception describe failures that require the model to represent and act on a
distinction between what it believes and what it reports. This design tests
neither, and we cite them to mark the boundary rather than to claim the
territory.

\citet{wen2024} is the closest prior result: models trained with human feedback
learn to mislead evaluators without being trained to, which they call
U-Sophistry. It is load-bearing here for a second reason. They report that a
probe trained to detect intended sophistry does not generalize to the
unintended case, and our pre-registration adopts that as the interpretation of an
E3 null. A near-chance probe result is a cross-task transfer failure, not
evidence that no fabrication signal is internally represented.

\citet{qi2026rewardseeker} approach the same problem from the training side.
They ran reinforcement learning on 80 environments they had identified as
vulnerable to reward hacking, and the resulting policy carried misaligned
behavior far outside that training distribution. Two of their observations bear
directly on this paper. They track a family of reward-seeking reasoning
patterns, ``related to how the model is being graded and its desire to maximize
its score on the episode,'' and one of them, grader sycophancy, was flagged in
27\% of episodes. That is the internal counterpart of the behavioral effect our
E2 measures from the outside. And their audit finding parallels our E4: ``Across
our production automated behavioral audit and many of our standard alignment
evaluations, the model overall appears as aligned as the initial checkpoint.''
Two designs in unrelated domains report that the standard evaluation apparatus
does not see the failure it exists to catch.

The boundary needs marking in three places, and the contrast is the useful part.
Their regime is reinforcement learning on task environments, while ours is the
supervised, preference and instruct stages of ordinary post-training, so their
result does not establish our cause. Their model is frontier-scale and ours are
8B to 72B open weights, the same tier mismatch we flag for
\citet{solutionhacking2026}. And they characterize their model as a
``reward-on-the-episode seeker'' that ``expresses motivation to achieve high
reward,'' which is a claim about what the model represents; our design observes
behavior only and we do not adopt that vocabulary. Set against that, they
induced the behavior deliberately in environments hand-selected for
exploitability, whereas we find a milder form already present in released
checkpoints that nobody set out to corrupt. Their subject is not citation
behavior, which the post does not discuss.

\subsection{Concurrent work on shortcut hacking}

\citet{solutionhacking2026} study what they call solution hacking on frontier
science benchmarks: a model reaches the correct answer through an invalid
shortcut such as enumeration or answer-first verification, bypassing the
reasoning the benchmark meant to measure. It is concurrent work, posted while
our data collection was underway, and it is a preprint rather than a
peer-reviewed result.

The constructs separate cleanly, which forecloses any priority question. Standard
benchmark scoring asks whether the answer is right and detects wrong answers.
Their audit asks whether a right answer was reached legitimately and detects a
correct outcome reached by an invalid process. Ours asks whether a cited artifact
has a referent at all and detects something that looks right and refers to
nothing.

Their construct boundary is useful to us for a structural reason. Their Table~1
distinguishes legitimate from illegitimate reasoning by the role a step plays in
a derivation rather than by the technique used, so the same enumeration can fall
on either side. Our §6(b) dividing line has the same shape: an identical
metadata discrepancy is an ordinary error when it points at the same work and a
fabrication when it identifies a different real one. That part of our taxonomy is
the most exposed to the objection that it is arbitrary, and it helps to be able
to point at an independent construct built the same way.

The empirical convergence is the stronger connection. Their hack ratio rises
monotonically with difficulty and does so within every subject, from 2.2\% on
their easy tier (SciBench and MATH-500) to 28.3\% on their medium tier
(IMO-Bench, PHYBench and SciOlympiad) to 37.4\% on their hard tier (HLE). Two
designs in different domains, measuring different failures, find the same
gradient: the behavior scales with how hard legitimate success is. They also
report that weaker models hack more, which is a useful external anchor for
reading the E1 cascade as a statement about difficulty rather than about a
particular model generation.

Their four anti-hack prompt variants all permit abstention, justified as our
\texttt{abstention\_permitted} condition is: without an abstention option the
prompt pushes a model to conceal a shortcut rather than admit it cannot solve
the problem. Their abstention rate moved from 0\% to 27.3\%. Ours is near zero
and quasi-separated. That divergence should not be over-read. Their population
is frontier 2026 models under an elaborate multi-variant intervention and ours
is 8B to 72B open-weights checkpoints under a single condition line, so the
honest reading is that the abstention lever is real and our near-zero rate is a
property of the model tier we studied.

On instrument quality the comparison cuts both ways, and we state it in both
directions. Their validation is stronger than ours: 300 solutions labeled by blinded
PhD annotators, 8, 7 and 6 per subject, with released guidelines, a held-out
test split used once, and a deployed detector that is a majority vote of three
peer-level judges under a no-self-audit rule. Our labels come from a single
annotator, so we have no inter-annotator reliability and no Cohen's $\kappa$.
That is the gap to concede, and it is a different gap from having an unvalidated
judge, which is not our situation (Section~\ref{sec:instrument}). Against that,
they report no temporal-stability or instrument-drift analysis, and their judge
is a static API call that cannot degrade the way a search-using judge can.

We also borrow their reporting standard, with the sign reversed. They report
that their detector misses 39\% of expert-confirmed hacks, and therefore frame
every rate they publish as a lower bound. Our judge fails in the opposite
direction under degraded grounding: it fails to confirm works that are real,
pushing verdicts toward fabrication, so our degraded-era rates are upper bounds.
Telling a reader which way to discount is more useful than reporting a single
accuracy number.

\subsection{Measuring reward hacking by construction, and auditing passes after the fact}

% ARGUES: our audit-based design is a considered choice forced by the object,
%         not an oversight, and the audit-and-overturn move has prior art.
% STATUS: drafted 2026-08-27. Citation metadata for the three new entries was
%         confirmed at the arXiv abstract pages. The quantitative figures below
%         (219 flaws, 16\%) come from those abstracts and still need a full-text
%         read before the first external compile.

Two recent lines of work bear on whether a claim like ours can be measured at
all, and we position against both rather than around them.

\citet{roth2026hackverifiable} argue that studying reward hacking by inspecting
trajectories after the fact is unreliable, and propose the alternative of
embedding detectable hacking opportunities directly into environments so that
exploitation becomes deterministically verifiable. We cannot take that route,
and the reason is a property of the object rather than a shortfall of effort. A
citation hack cannot be planted, because the opportunity is intrinsic: any model
asked to support a claim can emit a well-formed reference, and nothing has to be
placed in the environment for it to do so. Nor can the check be made local.
Deciding whether a reference has a referent is an open-world lookup against the
published literature, not a predicate evaluable inside a sandbox. The
obscure-but-real cell is the closest substitute the domain admits, since it
holds capability roughly fixed while changing what an honest answer looks like,
but it plants no exploit and yields no deterministic verifier.

\citet{qi2026rewardseeker} take the by-construction route one level up, building
the vulnerable training environments rather than the vulnerable task and reading
the disposition that results. That is available at the training level precisely
because the exploit can be placed there, which is the move the citation object
denies us at the task level.

The consequence is that our detection is an audit, and audits have an error rate
that a constructed verifier does not. We treat that as a cost to be measured
rather than assumed away, which is what Section~\ref{sec:instrument} reports.

That same audit-and-overturn structure has precedent, in a different domain.
\citet{wang2026benchjack} red-team ten agent benchmarks and surface 219 distinct
flaws across eight recurring classes, and build working exploits that pass tasks
without solving them. \citet{zhong2026hardening} audit 1,968 tasks across five
benchmarks with hand-written
verifiers and find 16\% exploitable, then close the loop by hardening the
verifiers themselves. Both establish the pattern we also rely on: a pipeline's
own passing verdicts are evidence about the pipeline, and re-checking them
recovers a non-trivial fraction that should not have passed.

The differences are what make our instance a separate contribution rather than
an application. The object differs, since theirs is a benchmark task with a
grading harness and ours is a generated citation with no harness at all. The
mechanism differs, since theirs is automated red-teaming against a verifier we
could inspect and ours is a language-model judge doing live web lookups, which
introduces a failure mode their verifiers do not have and cannot have: the
instrument's own grounding can degrade over time. And the remedy differs. They
patch the verifier so the exploit stops working; we cannot patch a model we did
not train, so we re-judge and report the correction, including its direction.

Industry practice has converged on the same move. The Artificial Analysis Coding
Agent Index began, in version 1.4 of August 2026, zeroing Terminal-Bench
attempts found to be reward hacking, stating that ``Flagged attempts are scored
zero'' and counting among its criteria the retrieval of a reference solution
from an external source by web search or fetch
\citep{artificialanalysis2026codingagents}. Treating a solution found on the open
web as a hack is a contamination argument as much as a reward-hacking one, and it
applies directly to the release of our own question set, which we take up in
Section~\ref{sec:limitations}. They publish no rates, so we cite the practice
and not a number.
